%% kapa: exam=2 round=22 year=2011 period=2 subject=appraisal_theory subject_ko=감정평가이론 seq=2 number=2 points=30
%% source: Q-Net 감정평가사 기출문제 (cst00309) artlSeq=5040389 file=extracted/1008/5040389/2011년 제22회 감정평가사 2차시험 기출문제/2교시.hwp
%% conversion: kordoc 4.24.51 -> tools/gen2.py (auto); review=proofread-sonnet; flags=none
\begin{longquestion}{【문제 2】}
부동산 감정평가 시 다양한 평가방법이 있고 정확한 가격 평가를 위해서는 경제적 상황의 변화도 고려해야할 필요가 있다. 다음의 물음에 답하시오. (30점) \par
\medskip\noindent (1) 감정평가에 사용될 수 있는 계량적(정량적) 방법인 특성가격함수모형(Hedonic Pricing Model)에 대해 설명하고, 감정평가사의 주관적 평가와 비교하여 그 장ㆍ단점을 논하시오. (10점) \par
\medskip\noindent (2) 최근의 세계경제 위기가 국내 부동산시장에 미치는 영향을 기술하고, 이러한 영향 하에서 부동산 감정평가를 할 경우 비교방식, 원가방식, 수익방식별로 유의점을 논하시오. (20점) \par
\end{longquestion}
